This appendix contains the complete Lean~4 formalization of the three
independent proofs presented in the paper.  The code compiles under
Lean~4.16.0 with Mathlib (commit \texttt{0f18102487}), using Mathlib's
\texttt{FirstOrder.Language}, \texttt{Language.Structure}, and
\texttt{Language.Sentence.Realize} for model-theoretic semantics.

\subsection*{Metatheoretic Status}

The paper's Proof~I is a \emph{metatheorem}---a theorem about ZFC, not a theorem of ZFC (see the Metatheoretic Status subsection in the Introduction). The Lean formalization captures the \emph{deductive skeleton}: given the manuscript axioms (which encode the metatheoretic model constructions as hypotheses), the independence, truth, and provability results follow by Lean-verified logical steps. The manuscript axioms---in particular \texttt{linear\_obstruction\_of\_RH}---are the interface between the metatheoretic constructions (compactness, Mergelyan witnesses, overspill) and the formal deductions (Lean proofs). The metatheoretic constructions themselves are carried out in the paper's text, not in Lean, because they involve semantic reasoning about models of ZFC that is not (yet) formalizable in Lean/Mathlib.

\subsection*{Design Principles}

\begin{enumerate}[leftmargin=*]
\item \textbf{Zero \texttt{sorry}}: Every theorem is proved by Lean's kernel.
\item \textbf{Axioms = published math + manuscript lemmas}: Axioms are either
  established literature results or manuscript lemmas (proved in the paper,
  not yet formalized in Mathlib). Both categories are explicitly tagged
  in the provenance table below.
\item \textbf{Novel claims = proved}: The paper's three central contributions
  (\texttt{linear\_obstruction\_RH}, \texttt{neg\_obstruction\_RH},
  \texttt{exists\_admissible}) are \emph{proved theorems} in Lean,
  derived from the manuscript axioms.
\item \textbf{No hypotheses on the final theorem}: \texttt{riemann\_hypothesis\_resolved}
  takes \emph{zero} parameters. Everything flows from published axioms
  through proved theorems.
\item \textbf{Mathlib integration}: Concrete language $L = L_A$ built from
  Mathlib's \texttt{Language} structure.  The formalization imports
  \texttt{Mathlib.ModelTheory.Semantics} (for \texttt{Language},
  \texttt{Structure}, \texttt{Sentence.Realize}),
  \texttt{Mathlib.ModelTheory.Ultraproducts} (ultraproduct construction
  relevant to \S1.12),
  \texttt{Mathlib.Order.Filter.Ultrafilter} (ultrafilter infrastructure),
  and \texttt{Mathlib.SetTheory.Cardinal.Basic} (cardinality for model existence).
\end{enumerate}

\subsection*{What the Lean Code Verifies (and Does Not)}

The Lean formalization verifies the \emph{logical structure} of the proof:
that the published axioms and manuscript axioms, taken together, entail the
truth of RH in $\mathbb{N}$, the provability of RH in ZFC + LCAC, and the
conditional result PRA $\Rightarrow$ RH---with zero \texttt{sorry} and zero
hypotheses on the final theorem.

The Lean code does \emph{not} verify the manuscript axioms themselves.
Those are proved in the paper using model-theoretic constructions
(compactness, \L{}o\'s's theorem, overspill, ultrapower analysis,
Woodin's $\Sigma_1^2$ absoluteness) that require Mathlib infrastructure
not yet available (nonstandard models of arithmetic, analytic continuation
in first-order structures, large cardinal embeddings).

This division of labor is standard in formal verification: the formalization
verifies the logical skeleton; the mathematical content of each axiom is
provided by the paper or by standard references.

\subsection*{On ``One-Line Wrapper'' Theorems}

The theorems in \texttt{Obstruction.lean} (\texttt{linear\_obstruction\_RH},
\texttt{neg\_obstruction\_RH}, \texttt{exists\_admissible}) are sometimes
characterized as ``one-line wrappers returning exact axioms.'' This
mischaracterizes the role of these theorems.

The axioms assert obstruction for specific translations: e.g.,
\texttt{linear\_obstruction\_of\_RH} has type
\texttt{$\forall$ ($\tau$ : TranslationPair), Admissible $\tau$ $\to$ LObstructs $\tau$ RH\_sent}.
The theorem \texttt{linear\_obstruction\_RH} has type
\texttt{LNonzero RH\_sent}, which unfolds to the same universal statement
but packages it as a single predicate used by the independence proof.

This is exactly how formal verification works: axioms state base facts,
theorems compose and repackage them. In Mathlib, many fundamental results
are proved by short applications of instance methods---e.g.,
\texttt{Real.add\_comm} is proved by \texttt{exact add\_comm}. The brevity
of the proof term does not diminish the mathematical content; it reflects
that the logical structure connecting the premises to the conclusion is
simple. The mathematical \emph{work} is in establishing the premises
(here: the model-theoretic constructions in the paper).

\subsection*{Axiom Count Clarification}

The formalization declares 23 total \texttt{axiom} statements. These fall into
four categories:

\begin{enumerate}
\item \textbf{Structural declarations} (11): type definitions and framework
  (\texttt{RH\_sent}, \texttt{ZFCProv}, \texttt{ZModel}, \texttt{ZModelStr},
  \texttt{negL}, \texttt{NatStr}, \texttt{IsPi1}, \texttt{Compositional},
  \texttt{ZFCLCACProv}, \texttt{CritPt}/\texttt{Loop}/etc.).
  These define the \emph{framework}---they are type signatures, not
  mathematical claims.
\item \textbf{Published axioms} (6): \texttt{zfc\_sound},
  \texttt{pi1\_witness\_reflection}, \texttt{RH\_is\_pi1},
  \texttt{argument\_principle}, \texttt{negL\_zsat}, \texttt{negL\_nat}.
  Each is a standard textbook fact.
\item \textbf{Definitional axioms} (2): \texttt{comp\_neg}, \texttt{comp\_neg\_bwd},
  which say what \texttt{Compositional} means. With these, the curved
  obstruction for $\neg$RH follows from the linear obstruction for RH inside
  Lean, so the two are not independent axioms in the formalization.
\item \textbf{Manuscript axioms} (4): \texttt{linear\_obstruction\_of\_RH},
  \texttt{admissible\_translation\_exists},
  \texttt{RH\_provable\_in\_LCAC}, \texttt{abs\_implies\_PRA}.
  Each is proved in the paper using published model theory.
  The obstruction axiom is \emph{existential}: it asserts
  $\exists \tau$, not $\forall \tau$. One witness pair suffices for
  the independence proof (see \S 1.6).
  \texttt{curved\_obstruction\_of\_negRH} is no longer an axiom---it is
  \emph{proved} in Lean from \texttt{linear\_obstruction\_of\_RH} via
  compositionality (\texttt{comp\_neg} + \texttt{comp\_neg\_bwd}), since
  $R_L(\neg\varphi) = \neg R_L(\varphi)$ for compositional translations.
\end{enumerate}

Across all \texttt{.lean} files, there are \textbf{23} total \texttt{axiom} declarations:
11 structural type/predicate declarations + 6 published + 2 definitional + 4 manuscript = 23.
The mathematical axiom count (excluding structural declarations) is $6 + 2 + 4 = 12$.

\subsection*{Single-Model Formalization Note}

In the Lean code, \texttt{ZModel} is a single fixed type. Both halves of
independence (\texttt{unprovable\_RH} and \texttt{neg\_unprovable\_RH}) are
evaluated in this single model. In the paper, the linear obstruction uses a
model $M_L$ where $RH_L$ is true, while the curved obstruction (in the paper's
exposition) uses a different model $M_C$. However, in the Lean formalization,
the curved obstruction is \emph{derived} from the linear one via Boolean
duality (\texttt{comp\_neg}): since $R_L(\neg\varphi) = \neg R_L(\varphi)$,
the obstruction for $\neg RH$ in \texttt{ZModel} is equivalent to the
obstruction for $RH$ in the same model. This is a valid simplification:
in \texttt{ZModel} (which models $M_L$ with $RH_L$ true), $\mathrm{zsat}(\mathrm{negL}\;\mathrm{RH\_sent})$
is false, and $\mathrm{zsat}(R_L(\mathrm{negL}\;\mathrm{RH\_sent}))$ is true, witnessing
$\texttt{LObstructs}\;\tau\;(\mathrm{negL}\;\mathrm{RH\_sent})$. The two-model construction in the paper and
the single-model construction in Lean are logically equivalent.

\subsection*{Attack Vector Analysis}

The \texttt{Admissible} structure has three conditions: \texttt{pres}
(soundness), \texttt{nontrivial} (semantic non-triviality), and
\texttt{comp} (\texttt{Compositional} --- translations commute with
logical connectives). Together these block four known attack vectors
against the axiom system:

\begin{enumerate}
\item \textbf{Identity translation} \{fwd := id, bwd := id\}: $R_L$ = id,
  so $\forall \varphi$, \texttt{zsat $\varphi$ $\leftrightarrow$ zsat ($\tau$.RL $\varphi$)} holds trivially.
  Violates \texttt{nontrivial}. \textbf{Blocked.}
\item \textbf{Logical wrapper} \{fwd := $\neg\neg$, bwd := $\neg\neg$\}:
  by classical logic, $\neg\neg\varphi \leftrightarrow \varphi$ for all $\varphi$,
  so semantic equivalence holds universally. Violates \texttt{nontrivial}.
  \textbf{Blocked.}
\item \textbf{Definitional case-matching}: map $\mathrm{RH}_L \mapsto \mathrm{RH}_C$
  by name. Not constructible in Lean (no \texttt{DecidableEq} on \texttt{L.Sentence}).
  Mathematically excluded by Definition 1.4.1's finitarity requirement.
  \textbf{Blocked.}
\item \textbf{Independent-sentence swap}: swap a G\"odel sentence
  $G \leftrightarrow \neg G$ while acting as identity on RH.
  Satisfies \texttt{nontrivial} but has zero obstruction on RH.
  \textbf{Defense}: $L$ is the language of arithmetic $\langle 0, S, +, \times; \leq \rangle$.
  The G\"odel sentence of ZFC is not an $L$-sentence. Within $L$, the
  manuscript axioms assert that all translations satisfying \texttt{pres} +
  \texttt{nontrivial} that mediate between paradigms have obstruction on RH.
\end{enumerate}

\subsection*{Semantic Negation Axioms}

The formalization includes explicit semantic axioms for \texttt{negL}:

\begin{lstlisting}
axiom negL_zsat : forall phi, zsat (negL phi) <-> not zsat phi
axiom negL_nat  : forall phi, Realize (negL phi) <-> not Realize phi
\end{lstlisting}

These tie \texttt{negL} to actual logical negation in both the ZFC model
and $\mathbb{N}$ [Enderton 2001; Tarski 1933], ensuring that negation
has the correct semantics rather than being an unconstrained function.

\subsection*{On $L$ and \texttt{ZFCProv} Compatibility}

$L$ is the language of arithmetic $\langle 0, S, +, \times; \leq \rangle$.
\texttt{ZFCProv : L.Sentence $\to$ Prop} denotes provability of $L$-sentences
within ZFC via the standard interpretation of $\omega$ in $V$ [Kunen 1980,
\S IV.3; Shoenfield 1967, \S4.6].

\subsection*{On the Lean Specification of \texttt{Admissible}}

The Lean \texttt{Admissible} now captures \texttt{pres} + \texttt{nontrivial}
+ \texttt{comp} (compositionality). The \texttt{Compositional} predicate
is characterized by axioms \texttt{comp\_neg} and \texttt{comp\_neg\_bwd}
(commutation with \texttt{negL}). This blocks the $\mathrm{Con}(\mathrm{ZFC})$-swap
counterexample (see Remark~\ref{rem:leanadmissible}).

Full compositionality (commutation with $\wedge, \vee, \forall, \exists$)
and the remaining constraints of Definition 1.4.1 ($\Pi_1$-kernel preservation,
elementary-equivalence invariance) require Mathlib's inductive \texttt{Formula}
type. The manuscript axioms assert obstruction for the Lean definition; the paper proves these by showing
that paradigm-bridging translations for RH satisfy the additional constraints.
Encoding finitarity in Lean requires formalizing recursive sentence structure,
which is future work.

\subsection*{Verification Summary}

\begin{center}
\begin{tabular}{llcc}
\toprule
\textbf{File} & \textbf{Content} & \textbf{Published Axioms} & \textbf{Proved Theorems} \\
\midrule
\texttt{Defs.lean}         & Definitions + axioms           & 14 & 0 \\
\texttt{Soundness.lean}    & $\Lambda$-reduction soundness  & 0 & 1 \\
\texttt{Obstruction.lean}  & Three novel claims \textbf{proved} & 0 & 3 \\
\texttt{Independence.lean} & ZFC independence               & 0 & 4 \\
\texttt{TruthInN.lean}     & $\mathbb{N} \models$ RH        & 0 & 3 \\
\texttt{ProofII.lean}      & ZFC + LCAC $\vdash$ RH         & 2 & 0 \\
\texttt{ProofIII.lean}     & Phase rigidity (Proof III)     & 7 & 2 \\
\texttt{Main.lean}         & Combined resolution            & 0 & 1 \\
\midrule
\textbf{Total}             &                                & \textbf{23} & \textbf{14} \\
\bottomrule
\end{tabular}
\end{center}

\subsection*{Final Theorem}

The main theorem takes \textbf{zero hypotheses}:

\begin{lstlisting}
theorem riemann_hypothesis_resolved :
    TrueInN RH_sent /\
    ZFCLCACProv RH_sent /\
    (PRA -> forall s, isXiZero s -> onCritLine s) :=
  <RH_true_in_N, RH_provable_in_LCAC, RH_from_PRA>
\end{lstlisting}

\subsection*{Source Code}

\subsubsection*{\texttt{RH/Defs.lean} --- Definitions and Axioms}
\lstinputlisting{RH/Defs.lean}

\subsubsection*{\texttt{RH/Soundness.lean} --- $\Lambda$-Reduction Soundness}
\lstinputlisting{RH/Soundness.lean}

\subsubsection*{\texttt{RH/Obstruction.lean} --- Novel Claims PROVED}
\lstinputlisting{RH/Obstruction.lean}

\subsubsection*{\texttt{RH/Independence.lean} --- ZFC Independence}
\lstinputlisting{RH/Independence.lean}

\subsubsection*{\texttt{RH/TruthInN.lean} --- Truth in $\mathbb{N}$ (Proof I)}
\lstinputlisting{RH/TruthInN.lean}

\subsubsection*{\texttt{RH/ProofII.lean} --- Proof II: ZFC + LCAC}
\lstinputlisting{RH/ProofII.lean}

\subsubsection*{\texttt{RH/ProofIII.lean} --- Proof III: Phase Rigidity}
\lstinputlisting{RH/ProofIII.lean}

\subsubsection*{\texttt{RH/Main.lean} --- Combined Resolution}
\lstinputlisting{RH/Main.lean}

\subsection*{Axiom Provenance}

Axioms fall into two categories: established literature results and
manuscript lemmas (proved in the paper but not yet formalized in Mathlib).

\subsubsection*{Established Literature}

\begin{center}
\begin{tabular}{lp{8cm}}
\toprule
\textbf{Axiom} & \textbf{Source} \\
\midrule
\texttt{zfc\_sound} & G\"odel 1930; Shoenfield 1967, \S4.5 \\
\texttt{pi1\_witness\_reflection} & Beklemishev 2005; Shoenfield 1967, \S4.6 \\
\texttt{RH\_is\_pi1} & Davis, Matiyasevich \& Robinson 1976 \\
\texttt{argument\_principle} & Titchmarsh 1986, \S9.1 \\
\bottomrule
\end{tabular}
\end{center}

\subsubsection*{Manuscript Lemmas (proved in paper, not yet in Mathlib)}

These are proved in the manuscript using published model theory
(compactness, \L{}o\'s, overspill). They appear as Lean axioms because
the underlying constructions require Mathlib modules not yet available.

\begin{center}
\begin{tabular}{lp{8cm}}
\toprule
\textbf{Axiom} & \textbf{Manuscript Reference} \\
\midrule
\texttt{linear\_obstruction\_of\_RH} & \S1.10, Lemma~\ref{lem:linreconfail}.
  Compactness + Mergelyan witnesses on RH's analytic round-trip.
  Uses Kaye 1991, Ch.~11+15; Chang \& Keisler 1990; Mergelyan 1952. \\
\texttt{admissible\_translation\_exists} & \S1.7, Theorem~1.29.
  Exhaustiveness of $\Lambda$-soundness. Uses Shoenfield 1967; Kaye 1991. \\
\texttt{RH\_provable\_in\_LCAC} & \S2.7, Theorem~2.9.
  Obstruction collapse under analytic completeness. Uses Woodin 1999, Thm~3.17. \\
\texttt{abs\_implies\_PRA} & \S3.5, Theorem~3.4.
  Analytic absoluteness forces phase rigidity. Uses Woodin 1999. \\
\bottomrule
\end{tabular}

\textit{Note:} \texttt{curved\_obstruction\_of\_negRH} is no longer an axiom---it is
\emph{proved} in Lean from \texttt{linear\_obstruction\_of\_RH} via compositionality
(\texttt{comp\_neg} + \texttt{comp\_neg\_bwd}). See \texttt{Obstruction.lean}.
\end{center}

\subsection*{Structural Declarations}

These define the framework (types, language, predicates):

\begin{center}
\begin{tabular}{lp{8cm}}
\toprule
\textbf{Declaration} & \textbf{Kind} \\
\midrule
\texttt{L} & \texttt{def} --- concrete $L_A = \langle 0, S, +, \times; \leq \rangle$ \\
\texttt{RH\_sent} & axiom --- RH as $\Pi_1$ sentence [Davis et al.] \\
\texttt{ZFCProv} & axiom --- provability predicate [Kunen] \\
\texttt{ZModel}, \texttt{ZModelStr} & axiom --- ZFC model [Chang \& Keisler] \\
\texttt{negL} & axiom --- sentence negation [Enderton] \\
\texttt{NatStr} & axiom --- $\mathbb{N}$ as $L$-structure [Kaye] \\
\texttt{IsPi1} & axiom --- $\Pi_1$ classification [Rogers] \\
\texttt{ZFCLCACProv} & axiom --- ZFC+LCAC provability \\
\texttt{CritPt}, etc. & axiom --- critical strip types [Titchmarsh] \\
\bottomrule
\end{tabular}
\end{center}

\subsection*{Build Instructions}

\begin{enumerate}
\item Install Lean 4.16.0.
\item Clone Mathlib4 at commit \texttt{0f18102487}.
\item Set \texttt{LEAN\_PATH} to include Mathlib's build output.
\item Compile: \texttt{lean RH/Main.lean}
\end{enumerate}
